%% ---------------------------------------------------------------------
\section{Buying the comparator more compute}\label{sec:budget}

\looseness=-1 Cost normalization makes the comparisons in Sec.~\ref{sec:law}
interpretable, and Sec.~\ref{sec:parity} states its price: a comparator held
to its published budget inside our recipe may be undertrained, not weak, and
the two are not separable by construction. We measured that reservation rather
than leave it as a limitation, re-running both self-supervised comparators at
four times their pre-training budget and at every fraction from 1 to 25\%,
since every other effect here moves with the data regime.
Table~\ref{tab:budget} reports it, and
Secs.~\ref{sec:budget:simsiam}--\ref{sec:budget:features} work through it.
The objection is confirmed outright: it costs us an accuracy claim on
convolutional backbones, attaches a budget and a data regime to the claim that
survives on attention, and replaces one account of our two populations'
difference with another.

\begin{table}[pos=htbp]
\caption{Buying the comparators more compute. CIFAR-100, ResNet-18, gain
over the baseline by data fraction; three seeds per pre-trained cell, ten
for the baseline and the prior at 1, 5 and 10\%. SSm is SimSiam and SC is
SimCLR, at 200 pre-training epochs ($2\times$ baseline compute) or 800
($5\times$); the prior costs $1.02\times$. The last two columns give the
prior minus each $5\times$ comparator, where positive
means the free prior still wins.}
\label{tab:budget}
\centering\scriptsize
\setlength{\tabcolsep}{2.6pt}
\zebra{6}
\begin{tabular}{lrrrrrrr}
\toprule
 & \textbf{Prior} & \multicolumn{2}{c}{$2\times$ SSL} & \multicolumn{2}{c}{$5\times$ SSL} & \multicolumn{2}{c}{prior $-$ $5\times$} \\
\cmidrule(lr){3-4}\cmidrule(lr){5-6}\cmidrule(l){7-8}
Frac. & $\mathbf{1.02\times}$~\up & SSm~\up & SC~\up & SSm~\up & SC~\up & SSm & SC \\
\midrule
1\%  & $+1.42$ & $-0.02$ & $+2.27$ & $-0.09$ & $\mathbf{+4.54}$ & $+1.51$ & $-3.12$ \\
2\%  & $+2.50$ & $-0.11$ & $+4.88$ & $+0.19$ & $\mathbf{+8.33}$ & $+2.31$ & $-5.83$ \\
3\%  & $+3.68$ & $+0.11$ & $+5.99$ & $+0.54$ & $\mathbf{+10.82}$ & $+3.14$ & $-7.14$ \\
5\%  & $+5.15$ & $+0.13$ & $+9.05$ & $+1.90$ & $\mathbf{+14.38}$ & $+3.25$ & $-9.23$ \\
7\%  & $+4.87$ & $+0.87$ & $+9.86$ & $+3.39$ & $\mathbf{+13.88}$ & $+1.48$ & $-9.01$ \\
10\% & $+3.75$ & $+0.61$ & $+8.76$ & $\mathbf{+4.33}$ & $\mathbf{+10.92}$ & $-0.58$ & $-7.17$ \\
15\% & $+2.55$ & $+0.17$ & $+6.64$ & $+3.65$ & $\mathbf{+8.36}$ & $-1.10$ & $-5.81$ \\
25\% & $+0.16$ & $-0.03$ & $+2.81$ & $+0.79$ & $\mathbf{+2.77}$ & $-0.63$ & $-2.61$ \\
\bottomrule
\end{tabular}
\end{table}

\subsection{The near-null was a budget artifact}\label{sec:budget:simsiam}
SimSiam at its published 200-epoch budget gains nothing anywhere on the
envelope, $-0.11$ to $+0.87$. At 800 epochs it gains up to $+4.33$. We
therefore withdraw the statement that negative-free self-supervision learns
almost nothing at this scale: it is budget-starved at the cost we normalized
to, and its published default is a poor guide to what it can do on a few
thousand images. Its budget response is itself data-dependent, \looseness=-1
buying most in the mid band ($+3.7$ at 10\%) and nothing at 1\%.

That data-dependence is why the envelope was necessary. The full
envelope shows a crossover: the prior wins from 1 to 7\% by up to $3.25$ ($2.6$
to $12.2$ standard errors) and \emph{loses} from 10 to 25\%, most clearly at
15\% ($-1.10$, $5.2$). A read at 5 and 10\% alone, where the objection is
sharpest, would have missed a significant
reversal, the same sparse-grid error we identify in our own earlier readings.

\subsection{On convolutions the prior's case becomes a cost
case}\label{sec:budget:conv} Against the stronger comparator the prior loses at
every fraction, including the scarcest. SimCLR at $5\times$ beats it by $3.12$
to $9.23$ points, at 12 to 37 standard errors, from 500 images to 7{,}500
(1--15\%), and by a narrowing $2.61$ at 25\% as both interventions decay toward
sufficiency. Our earlier positioning noted near-parity with self-supervision at
1\%; that was a statement at $2\times$ and it does not survive $5\times$. On
convolutional backbones the prior's case is therefore a cost case and not an
accuracy case, and we say so without qualification.

\subsection{Attention, given the same test}\label{sec:budget:vit}
Leaving the attention comparison at $2\times$ while raising the convolutional
one to $5\times$ would reproduce exactly the asymmetry this section exists to
remove, so we ran it there too. Against ViT-tiny under the DeiT recipe,
$5\times$ SimCLR gives $30.45$, $40.57$ and $55.85$ at 5, 10 and 25\%, against
the prior's $28.87$, $41.29$ and $57.32$. The comparator \emph{wins} at 5\% by
$1.57$ ($3.9$ standard errors, six seeds), the two are level at 10\% ($+0.71$,
$1.2$), and the prior wins at 25\% by $1.47$ ($7.2$). The prior is not
overturned on attention as it is on convolutions, but the claim must name its
budget: at $2\times$ the prior leads at every CIFAR-100 fraction, and at
$5\times$ only in the higher-data band. Section~\ref{sec:budget:tin} shows the same boundary
on a second population.

That shape contradicted a prediction recorded before the runs: we expected the
prior to hold at the scarce fractions and 25\% to be the cell most likely to
flip, and the reverse happened. What the extra budget buys shows why: $+8.80$,
$+4.51$ and $+3.53$ at 5, 10 and 25\%, decreasing with data, because 200 epochs
over 2{,}500 images is fewer than four thousand steps: the comparator was
step-starved at 5\%, not data-starved (Sec.~\ref{sec:dataopt}). On
convolutional backbones the same increment peaks in the mid band, so the budget
response is itself backbone-dependent.

\subsection{A second population, and the account that
failed}\label{sec:budget:tin} Repeating the comparison on Tiny-ImageNet gives
the same ordering at every matched fraction. Writing the prior minus a
$5\times$-budget SimCLR initialization, both under DeiT augmentation on
ViT-tiny, CIFAR-100 gives $-1.57$, $+0.71$ and $+1.47$ at 5, 10 and 25\%, and
Tiny-ImageNet gives $-1.35$, $-0.31$ and $+1.71$. The comparator wins at 5\% on
both ($3.9$ and $3.7$ standard errors), the two are level at 10\% on both, and
the prior wins at 25\% on both ($7.2$ and $5.1$). So at five times the
comparator's pre-training budget the prior leads in the higher-data band and
loses in the scarcer one, and that pattern is not specific to a dataset. The
ordering tracks the data fraction, not the absolute image count: $5{,}000$
images is 10\% of CIFAR-100 but 5\% of Tiny-ImageNet, and the two disagree
there while agreeing at every matched fraction.

Getting there cost us a prediction and an account. With only the 5 and 10\%
cells measured, where the comparator was ahead and level, we wrote that the
prior led nowhere on this population, predicted the 25\% cell at $40$--$44$
explicitly against our own method, and expected the two populations to line
up by image count. The cell landed at $36.53$, below the band; the recorded
falsifier fired; and the matched-fraction table shows why the account was
wrong: matching on image count compares different points on two different
learning curves, and what determines which source wins is where a cell sits
on its own dataset's curve. The ``different stories'' reading was an
artifact of a two-point grid, the same error this section has already
documented once.

One measurement in this pass we report without explaining. At Tiny-ImageNet
25\%, four times the pre-training made the comparator \emph{worse}: $38.48$ at
$2\times$ against $36.53$ at $5\times$, a drop of $1.95 \pm 0.42$. It is the
only cell in the study where extra pre-training budget costs accuracy, and it
echoes our earlier finding that DeiT-strength contrastive views also hurt. We
note it and decline to fit an explanation to one cell.

\subsection{What the budget actually buys}\label{sec:budget:features}
The comparison so far is end-to-end, while the taxonomy is a claim about
\emph{what} a source supplies, not how much accuracy it yields, so we
probed these cells under the frozen-feature protocol used throughout.

On CIFAR-100 the feature gain of $5\times$ SimCLR is $+22.24$, $+21.27$ and
$+20.77$ at 5, 10 and 25\%, against the prior's $+21.37$, $+22.19$ and
$+23.05$. The two orderings agree cell by cell: the comparator's feature gain
is larger where it wins end-to-end, level where the two are level, and smaller
where the prior wins. The prior's feature gain \emph{rises} with data while the
comparator's \emph{falls}, and they cross between 5 and 10\%, which is where
the accuracy ordering crosses. Extra pre-training budget therefore buys more of
the same currency, not a different one, and the flip is a property of the
representations, not of the classifier reading them. This is the substitution
outcome measured on a budget axis, the one axis the taxonomy had not been
tested on.

\looseness=-1 Including the second population strengthens the agreement and supplies the
mechanism the image-count account lacked. Across both datasets and all six
measured fractions, whichever intervention has the larger frozen-feature gain
wins end-to-end: the sign of $G_{\mathrm{SSL}} - G_{\mathrm{prior}}$ matches
the sign of the accuracy difference six times out of six. The clearest case is
the cell that overturned our own prediction. At Tiny-ImageNet 25\%, where the
prior wins by $1.71$, its feature gain is $+14.27$ against the comparator's
$+12.43$; that difference of $1.84$ agrees with the accuracy difference to
within $0.13$, so the outcome is feature-side almost entirely. The population
term is equally unmysterious once measured. At a matched 5{,}000 images the two
datasets diverge because the \emph{prior's} contribution diverges: its feature
gain is $+22.19$ on CIFAR-100 against $+15.67$ on Tiny-ImageNet, while the
comparator's is $+21.27$ and $+17.72$. Moving to the 200-way $64$\,px task
costs contrastive pre-training about $3.5$ points of feature gain and the fixed
spectral target about $6.5$: the hand-crafted prior supplies proportionally
less of what a finer label space at higher resolution requires, and that,
not the number of images, decides which source wins.\footnote{The
Tiny-ImageNet comparator probes were computed on a different machine from the
cells they pair with. Cross-machine agreement under this protocol was verified
elsewhere in the study to within $0.05$ points, but we note the mixed
provenance instead of leaving it implicit.}

The second comparator extends the same agreement. Probing SimSiam at $5\times$
across all eight fractions gives a feature gain that sits below the prior's
from 1 to 7\% ($-5.3$ to $-1.9$ points of difference), crosses at 10\%, and
rises above it at 15\%, matching the end-to-end ordering reported in
Sec.~\ref{sec:budget:simsiam} at seven of eight cells; the one mismatch is
the crossing cell itself, where the feature difference is a near-zero $-0.08$. Across two
self-supervised methods and two backbone families, the source with the larger
frozen-feature gain is the one that wins.

\looseness=-1 One result in this pass we report without an account. Probing the convolutional
budget envelope shows the prior's own residual tracking the fitted branch
closely ($-2.74$, $-2.65$, $-2.07$, $-1.11$ below the crossing, $+0.06$ and
$+0.21$ inside the bracket, and $-0.18$, $-0.28$ above it), while $5\times$ SimCLR departs from it,
running positive from 3 to 15\% ($+1.32$, $+3.11$, $+3.58$, $+2.23$, $+1.81$)
at baselines where the branch is negative or near zero. Two things narrow what
the anomaly can be. It is not a property of self-supervised initialization as
such: $5\times$ SimSiam, initialized the same way and evaluated identically,
stays within $[-0.42, +0.86]$ at every fraction. And across all sixteen
self-supervised convolutional cells the residual correlates with the size of
the intervention itself ($r = +0.52$ between $\Delta$ and the residual), so
what departs from the branch is large gains rather than a particular method. We
recorded in advance that these initializations lie outside the scope in which
the branch was estimated, so we do not count the departure against it, and we
prefer to leave it as a measured regularity than to fit an explanation to
sixteen points.

\emph{Takeaway.} A comparator's budget is part of the claim. Two of our own
accuracy claims did not survive this section and are now cost claims; the
finding that did survive is that heavy augmentation, not the prior, is what
collapses contrastive pre-training's advantage.
